%iffalse
\let\negmedspace\undefined
\let\negthickspace\undefined
\documentclass[journal,12pt,twocolumn]{IEEEtran}
\usepackage{cite}
\usepackage{amsmath,amssymb,amsfonts,amsthm}
\usepackage{algorithmic}
\usepackage{graphicx}
\usepackage{textcomp}
\usepackage{xcolor}
\usepackage{txfonts}
\usepackage{listings}
\usepackage{enumitem}
\usepackage{mathtools}
\usepackage{gensymb}
\usepackage{comment}
\usepackage[breaklinks=true]{hyperref}
\usepackage{tkz-euclide} 
\usepackage{listings}
\usepackage{gvv}                                        
%\def\inputGnumericTable{}                                 
\usepackage[latin1]{inputenc}                                
\usepackage{color}                                            
\usepackage{array}                                            
\usepackage{longtable}                                       
\usepackage{calc} 
\usepackage{multirow}                                         
\usepackage{hhline}                                           
\usepackage{ifthen}                                           
\usepackage{lscape}
\usepackage{tabularx}
\usepackage{array}
\usepackage{float}


\newtheorem{theorem}{Theorem}[section]
\newtheorem{problem}{Problem}
\newtheorem{proposition}{Proposition}[section]
\newtheorem{lemma}{Lemma}[section]
\newtheorem{corollary}[theorem]{Corollary}
\newtheorem{example}{Example}[section]
\newtheorem{definition}[problem]{Definition}
\newcommand{\BEQA}{\begin{eqnarray}}
\newcommand{\EEQA}{\end{eqnarray}}
\newcommand{\define}{\stackrel{\triangle}{=}}
\theoremstyle{remark}
\newtheorem{rem}{Remark}

% Marks the beginning of the document
\begin{document}
\bibliographystyle{IEEEtran}
\vspace{3cm}
\title{JEE}
\author{AI24BTECH11004-BHERI SAI LIKITH REDDY}
\maketitle
\newpage
\bigskip

\renewcommand{\thefigure}{\theenumi}
\renewcommand{\thetable}{\theenumi}
\vspace{1cm}
\maketitle{Section-A JEE Advanced/ IIT-JEE}\\
\begin{enumerate}
\item $f(x)=\begin{vmatrix}
		$$\sec x$$ & $$\cos x$$ & $$\sec^2x+\cot x \cosec x$$\\
		$$\cos^2 x & \cos^2 x & \cosec^2 x$$\\
		$$1 & \cos^2 x & \cos^2 x$$
	\end{vmatrix}$\\
Then $\int_0^{\frac{\pi}{2}}$ $f(x) dx$ = .......
\hfill{(1987 - 2 Marks)}\\
		
\item The integral $\int_0^1.5 [x^2]dx$,
\hfill{(1988 - 2 Marks)}\\
Where []denotes the greatest integer finction, equals ......\\
		
\item The value of  $\int_{-2}^2 |1-x^2|dx$ is ...
	\hfill{(1989 - 2 Marks)}\\
		
\item The value of $\int_{\frac{\pi}{4}}^{\frac{3\pi}{4}} \frac{\phi}{1+ \sin \phi} d\phi$
\hfill{(1993 - 2 Marks)}\\
		
\item The value of $\int_2^3 \frac{\sqrt x }{\sqrt {5-x}+\sqrt {x}} dx$
\hfill{(1994 - 2 Marks)}\\
		
\item If for nonzero $x$, $af(x)+bf \left(\frac{1}{x} \right) = \frac{1}{x}-5 $ where $a$ 
	$\neq$  $b$, then $\int_1^2f(x)dx$=....
\hfill{(1996 - 1 Mark)}\\
		
\item If n>0, $\int_1^{2\pi} \frac {x \sin ^{2n}x}{\sin ^{2n} x + \cos ^{2n} x} dx $
\hfill{(1996 -1 Mark)}\\
		
\item The value of $\int_1^{e^{37}} \frac {\pi \sin (\pi lnx)}{x} dx$ is ...

\hfill{(1997 - 2 Marks)}\\
		
\item Let $\frac {d}{dx}F(x)=\frac {e^{\sin x}}{x},x>0$.If $\int_1^4 \frac{2e^{\sin x^2}}{x}=F(k)-F(1)$ then one of the possible valued of k is ...
\hfill{(1997 - 2 Marks)}\\
\end{enumerate}
\maketitle{Section B True/False}
\begin{enumerate}
\item The value of the intrgral $\int_0^{2a} \frac{f(x)}{[f(x)+f(2a-x)]}dx$ is equal to a
\hfill{(1988 - 1 Mark)}\\
\end{enumerate}
\maketitle{Section C MCQs with One Correct Answer}
\begin{enumerate}
	\item The value of the definite integral $\int_0^{2a}(1+e^{-x^2})dx$ is\\
		\begin{enumerate}
	\item -1            
	\item  2\hfill{(1981 - 2 Marks)}
	\item  1+$e^{-1}$     
	\item  none of these
		\end{enumerate}
	\item Let $a,b,c$ be non-zero real numbers such that $\int_0^1 (1+ \cos ^8 x)(ax^2 + bx +c)dx = \int_0^2 (1+ \cos^8 x)(ax^2+bx+c)dx.$\\
		Then the quadratic equation $ax^2+bx+c=0$ has
		\hfill{(1981 - 2 Marks)}\\
		\begin{enumerate}
	\item no roots in (0,2)
	\item at least one root in (0,2)
	\item double root in (0,2)  
	\item two imagenary roots
\end{enumerate}
	\item The area bounded by the curves $y=f(x)$, the x-axis and the ordinate x=1 and x=b is (b-1) sin (3b+4). Then $f(x)$ is 

		\hfill{(1982 - 2 Marks)}\\
		\begin{enumerate}
			\item$(x-1) \cos {3x+4}$
			\item$\sin {(3x+4)}$
			\item$\sin {(3x+4)}+3(x-1) \cos {(3x+4)}$
			\item none of the above
		\end{enumerate}
	\item the value of the integral $\int_0^{\frac{\pi}{2}} \frac{\sqrt{\cot x}}{\sqrt{\cot x}+\sqrt{\tan x}}dx$ is 

		\hfill{(1983 - 1 Mark)}\\
		\begin{enumerate}
			\item $\frac{\pi}{4}$
		\item$\frac{\pi}{2}$
			\item
			$\pi$
		\item none of the above
		\end{enumerate}
	\item For any integer $n$ the integral-\\
		$\int_0^{\pi} e^{\cos^2 x} \cos ^3 {(2n+1)}xdx$ has the value

		\hfill {(1985 - 2 Marks)}\\
		\begin{enumerate}
			\item	$\pi$
			\item1
			\item0
			\item none of these
		\end{enumerate}
\end{enumerate}

    

\end{document}
